\documentclass[pdflatex,sn-basic]{sn-jnl}

\usepackage{amsmath,amssymb}
\usepackage{bm}
\usepackage{graphicx}
\usepackage{booktabs}
\usepackage{array}
\usepackage{textcomp}
\usepackage{xcolor}
\usepackage{manyfoot}
\usepackage[section]{placeins}

\raggedbottom

\begin{document}

\title[Supplementary Information]{Supplementary Information for ``A
Conservative Darcy-Brinkman-Stokes Formulation for PHREEQC-Coupled Reactive
Transport with Cell-Based Surface Complexation''}

\author*[1,2]{\fnm{Hannah P.} \sur{Menke}}\email{H.Menke@hw.ac.uk}
\author[1,2]{\fnm{Florian} \sur{Doster}}
\author[1,2]{\fnm{Julien} \sur{Maes}}

\affil*[1]{\orgdiv{Institute for GeoEnergy Engineering}, \orgname{Heriot-Watt University}, \orgaddress{\city{Edinburgh}, \country{United Kingdom}}}
\affil[2]{\orgdiv{SET-IT (Subsurface Energy Transition \& Innovation Centre)}, \orgname{Heriot-Watt University}, \orgaddress{\city{Edinburgh}, \country{United Kingdom}}}

\maketitle

+\section{One-Dimensional Numerical Convergence}

The spatial and temporal sensitivity of the 1D retardation
benchmark was evaluated at $t=10$~s for both the surface-enabled and matched
aqueous-only cases. The mesh study used $N_x=40$, 80, 160, and 320 cells with
$\Delta t=0.015625$~s. The timestep study used $N_x=320$ with
$\Delta t=0.25$, 0.125, 0.0625, 0.03125, and 0.015625~s. All other physical and
chemical inputs were unchanged from the benchmark reported in the main text.

Table~\ref{tab:1dConvergence} reports the relative $E_1$ difference between
each concentration profile and the next finer result. For the mesh comparison,
the finer profile was conservatively averaged onto the coarser cells before
evaluation. The successive differences decrease monotonically for both cases.
The observed orders range from 0.83 to 0.96 for the surface-enabled mesh
sequence and from 0.95 to 0.96 for its timestep sequence; the corresponding
aqueous-only ranges are 0.76--0.99 and 0.94--1.13. These values are consistent
with the first-order Euler time integration, upwind advection, and sequential
transport--equilibrium coupling used here. Because these are pairwise numerical
differences, they do not depend on a semi-infinite reference profile.

\begin{table}[htbp]
    \centering
    \caption{Mesh and timestep refinement for the 1D retardation benchmark at
    $t=10$~s. Values are relative $E_1$ differences (\%) from the next finer
    result. A dash identifies the finest result in each sequence.}
    \label{tab:1dConvergence}
    \begin{tabular}{llrrr}
        \toprule
        Study & Resolution & $N_x$ & Surface-enabled & Aqueous-only \\
        \midrule
        Mesh & $\Delta x=1.000~\mu$m & 40  & 0.949 & 0.583 \\
             & $\Delta x=0.500~\mu$m & 80  & 0.532 & 0.344 \\
             & $\Delta x=0.250~\mu$m & 160 & 0.274 & 0.173 \\
             & $\Delta x=0.125~\mu$m & 320 & --    & --    \\
        \midrule
        Timestep & $\Delta t=0.250000$~s & 320 & 1.444 & 0.465 \\
                 & $\Delta t=0.125000$~s & 320 & 0.747 & 0.213 \\
                 & $\Delta t=0.062500$~s & 320 & 0.383 & 0.109 \\
                 & $\Delta t=0.031250$~s & 320 & 0.198 & 0.057 \\
                 & $\Delta t=0.015625$~s & 320 & --    & --    \\
        \bottomrule
    \end{tabular}
\end{table}

The $80$-cell, $\Delta t=0.25$~s discretisation used for the main validation
differs from the finest combined result
($N_x=320$, $\Delta t=0.015625$~s) by $3.01\%$ and $1.05\%$ for the
surface-enabled and aqueous-only profiles, respectively. The monotonic
pairwise differences and observed orders quantify the combined spatial,
temporal, and sequential-coupling errors in the reported discretisation.

\section{Reactive-Area Construction and Capacity Mapping}

For each ordered pair of segmented labels, the owner and other phase define
a pair mask through their union. In \texttt{occupancyGaussian}, the
Gaussian-filtered owner indicator is divided by the Gaussian-filtered pair
mask to form a normalised occupancy field. The $0.5$ level is extracted
within the pair mask.

In three dimensions, reconstructed triangle area is assigned to computational
cells and divided by bulk cell volume. Paper-local three-dimensional
deposition probes owner points at offsets from $0.2$ to $3.0$ voxels along
the triangle normal, followed by a nearest-owner search in the surrounding
voxels. If this search finds no owner voxel, the triangle centroid supplies
the mapping point.

In quasi-two-dimensional images, reconstructed contour length is multiplied
by the physical out-of-plane thickness before division by bulk cell volume.
The serial construction probes offsets from $0.2$ to $1.5$ voxels along the
contour normal and uses the contour-piece midpoint if the probes find no
owner point. Weighted redistribution assigns each contour piece to
owner-side first-shell contact cells, with weights based on distance and
contact support. Candidate weights are normalised for each piece.
Local quasi-two-dimensional construction uses contact-support-based mapping
and deposits the weighted shares assigned to rank-owned cells. Pieces with
empty support contribute zero area to the deposited field. Both dimensional
routes return m$^2$\,m$^{-3}_{\mathrm{bulk}}$. The field is constructed before
flow and held fixed throughout reactive transport.

For intrinsic area, $s_k$ is defined per unit volume of phase $k$
(m$^2$\,m$^{-3}_{k}$). The geometric phase fraction $f_k$ converts this
end-member density to a bulk-volume cell contribution:
$a_{\mathrm{surf}}^{\mathrm{intrinsic}}=\sum_k f_k s_k$.
A mass-specific measurement is converted to phase-volume density using the
corresponding mass per unit phase volume. A density already expressed per
bulk cell volume represents the cell contribution directly.

The \texttt{phreeqcMixture} builds used for the reported calculations supply
PhreeqcRM \citep{parkhurst2015phreeqcrm} with area per representative bulk
litre, obtained by dividing $a_{\mathrm{surf}}$ in
m$^2$\,m$^{-3}_{\mathrm{bulk}}$ by $10^3$. For surface master $k$, the site
density $N_k$ (mol\,m$^{-2}$) gives the surface-master inventory per bulk
volume as $N_k a_{\mathrm{surf}}$. These bulk-volume conversions apply
without division by $\theta$. The inventory units are
\[
\left[\frac{\mathrm{m}^2}{\mathrm{m}^3_{\mathrm{bulk}}}\right]
\left[\frac{\mathrm{mol}}{\mathrm{m}^2}\right]
=\left[\frac{\mathrm{mol}}{\mathrm{m}^3_{\mathrm{bulk}}}\right].
\]

The chemistry mask threshold is set by \texttt{epsChemistryMin} in
\texttt{constant/thermoPhysicalProperties}, and every calculation reported here
uses $\theta_{\mathrm{mask}}=10^{-3}$.

\section{Discrete PhreeqcRM Coupling}

The surface-complexation cases select \texttt{mixtureType phreeqcMixture} in
\texttt{constant/thermoPhysicalProperties}. The active mapping between
finite-volume cells and PhreeqcRM cells is established once and does not change
during the run. Reconstruction is needed only when a threshold crossing changes
the active-cell map or its topology, or when the fixed area capacity changes.

Each OpenFOAM process constructs PhreeqcRM with one worker thread. Domain
decomposition provides process-level parallelism, and \texttt{runCells()}
advances only the chemistry cells held by that rank. Each case is decomposed
according to the \texttt{NPX}, \texttt{NPY}, and \texttt{NPZ} entries in its
\texttt{system} directory. Every case uses its case-local
\texttt{constant/GeoChem.dat} database. The image-based cases use kinematic
viscosity $\nu = 10^{-6}$~m$^2$\,s$^{-1}$ and kinematic pressure drop
$\Delta p = 5\times10^{-5}$~m$^2$\,s$^{-2}$, and write output every 10~s for the 1D
benchmark, 2000~s for the fracture case, and 20~s for the Estaillades case.

\subsection{Checkpoint and Restart Lifecycle}

At an output time, a checkpoint-enabled run first invalidates the old
completion marker collectively, writes a new rank-local dump, writes the
finite-volume fields, and publishes a completion marker only after the field
write succeeds. A crash before publication therefore leaves the partially
updated fields and chemistry state without a marker, and restoration rejects
them. The PhreeqcRM dump and marker must be retained with the corresponding
finite-volume time directory; the \texttt{Surf\_*} fields alone do not
reconstruct the complete internal state.

\section{Short Two-Scale Reconstruction Sensitivity}

The two-scale calculation is a matched short-interval sensitivity test between
two DBS cases that differ only in the surface field. Both use a
$3600\times1200\times1$ grid with $5~\mu$m cells,
16 MPI ranks, $Pe=1$, $D=10^{-9}$~m$^2$\,s$^{-1}$, and prescribed Darcy velocity
$5.46297\times10^{-5}$~m\,s$^{-1}$. The surface field is the only input that
differs, and the other 495 input files were verified as identical.

Both short cases use their case-local \texttt{GeoChem.dat} together with the
same illustrative Li-selective sorbent chemistry. The surface reactions are
$\mathrm{Sor\_aH}+\mathrm{Li^+}\rightleftharpoons\mathrm{Sor\_aLi}+\mathrm{H^+}$
($\log K=-5.8$), and the corresponding Na$^+$, K$^+$, and Mg$^{2+}$ reactions
have $\log K=-7.2$, $-7.0$, and $-7.4$, respectively. Initially Li-free brine
contains 0.10~M Na, 0.01~M K, 0.01~M Mg, $10^{-12}$~M Li, and 0.13~M Cl at pH
7; the inflow changes Li to $10^{-3}$~M and Cl to 0.131~M. These parameters
apply to the short $Pe=1$ sensitivity test only.

Three common checkpoints at 0.0058517566, 0.11703513, and 0.99564208~s meet
the predefined material-response limits. The maximum absolute differences
are $0.000647\%$ in integrated area, $0.000703\%$ in aqueous Li inventory,
$0.2110\%$ in sorbed Li inventory, $0.00172\%$ in the aqueous profile-shape
diagnostic, $0.5147\%$ in the sorbed profile-shape diagnostic, and $0.03993\%$
in fixed-reference-area sorbed-response L2. The two profile-shape diagnostics
compare independently normalised profiles and are not the relative $E_1$
measure used for validation. Shifts in the aqueous and sorbed cumulative-inventory median positions
($x_{50}$) are zero. The native-cell
sorbed inventory-density L2 is $3.61$ to $4.49\%$, dominated by the deliberate
relocation of area, so the agreement is in integrated response rather than
pointwise. The last accepted checkpoint covers $0.489\%$ of the configured
$203.53$~s horizon, so the test establishes early material-response agreement.

\section{Reproducibility and Provenance Entry Points}

The paper directory contains self-contained case definitions and source-data
references. From the repository root, the provenance verifier is run with
\begin{verbatim}
./validate_accepted_nonkelvin_artifacts.sh
\end{verbatim}
This command runs the verifier tests and validates
\texttt{figures/accepted\_nonkelvin\_provenance.json}. The manifest binds the
reported 1D, fracture, and Estaillades figures to their source data, renderers,
sidecars, and evidence records by path, size, and SHA-256.

The case-level entry point for the 1D benchmark is
\texttt{runPhase8AnalyticalBenchmark.sh} in
\nolinkurl{cases/1d_linear_retardation}. The mesh and timestep sequences are
reproduced by \texttt{runConvergenceStudy.sh}; the resulting metrics and exact
runtime identities are retained in \texttt{convergence\_results.json} and
\texttt{convergence\_provenance.json}. The fracture and Estaillades cases are
run through the \texttt{runAll.sh} scripts in
\nolinkurl{cases/2d_4StraightFractures} and \nolinkurl{cases/3d_estaillades},
each of which also provides separate mesh, flow, reactive-transport,
processing, and cleanup scripts. The short
two-scale record is in the \texttt{results} directory of
\nolinkurl{cases/two_scale_micromodel_sorption_validation}; it records exact
checkpoints, thresholds, collection hashes, runtime identity, and analyser
identity. The paper sources and both PDFs are built with \texttt{build.sh}. The
flat Springer submission archive is generated transactionally with
\texttt{export\_submission.sh OUTPUT.zip}; its manifest and SHA-256 inventory
are included in the archive.


\backmatter
\bibliography{references}

\end{document}
